% =========================================================================
% CONSOLIDATED STRATEGIC PRIORITY CLUSTERS WITH METHODOLOGY MATRICES
% File: Consolidated_Priorities_Matrices_v1.tex
% Standalone compilable.
%
% Each cluster has a phase-as-rows x 4-dimension-as-columns transition
% matrix, consistent with the generalized Methodology_Table_v2.tex.
%
% Dimensions per phase:
%   Col 1: R&I Maturity            (ODP scale)
%   Col 2: Policy, Regulatory & Standards
%   Col 3: Operational Deployment & Adoption
%   Col 4: Market Uptake & Strategic Sovereignty (SSL scale)
%
% Rows per matrix:  Phase 1 / Gate / Phase 2 / Gate / Phase 3
%
% Color scheme reused from Methodology_Table_v2.tex for consistency.
% =========================================================================

\documentclass[10pt,a4paper]{article}
\usepackage[T1]{fontenc}
\usepackage{geometry}
\usepackage{xcolor}
\usepackage{colortbl}
\usepackage{array}
\usepackage{longtable}
\usepackage{pdflscape}
\usepackage{enumitem}
\usepackage{booktabs}
\usepackage{hyperref}

% Standalone mode: render citation markers without requiring a bibliography pass.
\renewcommand{\cite}[2][]{\textit{[ref]}}

\geometry{margin=1.4cm}

% ---- Row background colours (identical to Methodology_Table_v2) ----
\definecolor{gaterow}{RGB}{240,240,238}
\definecolor{phase1row}{RGB}{237,245,253}
\definecolor{phase2row}{RGB}{237,249,244}
\definecolor{phase3row}{RGB}{245,244,254}
\definecolor{colheader}{RGB}{245,245,243}
\definecolor{clusterheader}{RGB}{220,232,248}

% ---- SSL badge colours ----
\definecolor{sslredbg}{RGB}{252,235,235}
\definecolor{sslredtx}{RGB}{163,45,45}
\definecolor{sslamberbg}{RGB}{250,238,218}
\definecolor{sslambtx}{RGB}{133,79,11}
\definecolor{ssltransbg}{RGB}{230,241,251}
\definecolor{ssltrantx}{RGB}{24,95,165}
\definecolor{ssltealbg}{RGB}{225,245,238}
\definecolor{ssltealtx}{RGB}{15,110,86}

% ---- Reused macros ----
\newcommand{\sslbadge}[3]{%
  \parbox{\linewidth}{\raggedright
    \colorbox{#1}{\scriptsize\parbox{0.92\linewidth}{\raggedright
      \textbf{\textcolor{#2}{#3}}}}}}

\newcommand{\kpistrip}[1]{%
  \par\vspace{2pt}\noindent\rule{\linewidth}{0.3pt}%
  \par\noindent\textit{\scriptsize #1}}

\newcommand{\cellhead}[1]{\noindent\textbf{\small #1}\par\medskip}

\newcommand{\gatelbl}[1]{\scriptsize\textit{Gate~#1:}}

\newcommand{\subp}[1]{\textit{\scriptsize(#1)}}

\begin{document}

% =========================================================================
% PREAMBLE
% =========================================================================

\section{Consolidated Strategic Priority Matrices}

The three strategic priority clusters are each assessed across four analytical
dimensions per phase: (1) R\&I Maturity, (2) Policy/Regulatory/Standards,
(3) Operational Deployment \& Adoption, and (4) Market Uptake \& Strategic
Sovereignty. Phases are presented as rows; the four dimensions are presented
as columns. Gate rows between phases state the minimum verifiable conditions
required to declare a transition. Sub-priority labels in parentheses
(\textit{e.g.}, P1, P3) identify where a specific constituent priority
divergences from the cluster-level descriptor within a given cell.

\bigskip
\noindent The three clusters and their constituent sub-priorities are:
\begin{itemize}[noitemsep]
  \item \textbf{Cluster~1 --- Sovereign Trusted Infrastructure:}
        P1 (Hardware Roots of Trust), P2 (PQC Agility),
        P3 (Lightweight Crypto \& PhySec), P8 (Confidential Computing),
        P9 (Resilient-by-Design / Island Mode).
  \item \textbf{Cluster~2 --- Cognitive Security \& Autonomous Cyber Defence:}
        P6 (Cognitive SecOps / LLMSecOps), P10 (Proactive APT Management),
        AI-as-target dimension (CL1.2, CL1.3, CL4.7).
  \item \textbf{Cluster~3 --- Federated Operational Sovereignty:}
        P4 (Security Interoperability via Simpl),
        P5 (Supply Chain Verification),
        P7 (Pan-European NaaS Federation).
\end{itemize}

\newpage

\newpage



\newpage


% =========================================================================
% CONSOLIDATED PRIORITY MATRICES — v2
% Progression matrices restructured to align with Methodology Table v2.
% Each matrix analyses all three phases across four analytical dimensions:
%   (1) R&I Maturity (ODP vocabulary)
%   (2) Policy, Regulatory & Standards
%   (3) Operational Deployment & Adoption
%   (4) Market Uptake & Strategic Sovereignty (SSL vocabulary)
% Plus a Critical Enabler (gate condition) row per phase transition.
% Sub-priority ODP/SSL decomposition retained as a nested lower section.
%
% Integration: replaces the three \resizebox progression matrix tables
% in consolidated_priorities_v1.tex. The 7-dimensional analysis sections
% in that file remain unchanged.
%
% Required packages (add to master preamble if not already present):
%   \usepackage[table]{xcolor}
%   \usepackage{longtable}
%   \usepackage{pdflscape}
%   \usepackage{array}
%   \usepackage{colortbl}
% =========================================================================

% ---- Shared colour definitions (safe to declare multiple times) --------
\definecolor{dimRI}{RGB}{237,245,253}       % R&I rows    — blue tint
\definecolor{dimPol}{RGB}{237,249,244}      % Policy rows — green tint
\definecolor{dimOps}{RGB}{253,243,237}      % Ops rows    — orange tint
\definecolor{dimMkt}{RGB}{245,244,254}      % Market rows — purple tint
\definecolor{dimGate}{RGB}{240,240,238}     % Gate rows   — neutral grey
\definecolor{dimSub}{RGB}{248,248,246}      % Sub-priority rows
\definecolor{dimHdr}{RGB}{220,220,218}      % Section separator rows
\definecolor{clr1hdr}{RGB}{46,117,182}      % Cluster 1 header bar
\definecolor{clr2hdr}{RGB}{84,130,53}       % Cluster 2 header bar
\definecolor{clr3hdr}{RGB}{112,48,160}      % Cluster 3 header bar

% ---- Reuse badge and strip commands from Methodology_Table_v2.tex ------
% (Define here only if not already defined in the master document)
\providecommand{\sslbadge}[3]{%
  \parbox{\linewidth}{\raggedright
    \colorbox{#1}{\scriptsize\parbox{0.92\linewidth}{\raggedright
      \textbf{\textcolor{#2}{#3}}}}}}
\providecommand{\kpistrip}[1]{%
  \par\vspace{2pt}\noindent\rule{\linewidth}{0.2pt}%
  \par\noindent\textit{\scriptsize #1}}
\providecommand{\cellhead}[1]{\textbf{\small #1}\par\smallskip}

% =========================================================================
% CLUSTER 1 MATRIX
% =========================================================================
\newpage
\section{Progression Matrix: v2 Format}

\begin{landscape}
\setlength{\tabcolsep}{4pt}
\renewcommand{\arraystretch}{1.25}

\begin{longtable}{%
|>{\footnotesize\raggedright\arraybackslash}p{3.4cm}%
|>{\footnotesize\raggedright\arraybackslash}p{5.5cm}%
|>{\footnotesize\raggedright\arraybackslash}p{5.5cm}%
|>{\footnotesize\raggedright\arraybackslash}p{5.5cm}|}

\caption{Progression Matrix — Priority Cluster~1: Sovereign Trusted
  Infrastructure (P1: Hardware RoT, P2: PQC Agility, P3: Lightweight
  Crypto \& PhySec, P8: Confidential Computing, P9: Island Mode).
  Each phase is assessed across four analytical dimensions per the
  generalised methodology framework.
  Critical Enabler rows state the single named gate condition controlling
  the transition into the following phase.}
\label{tab:cluster1_matrix}\\
\hline
\rowcolor{dimHdr}
\multicolumn{4}{|>{\footnotesize\bfseries\centering\arraybackslash}p{21.5cm}|}{%
  \textcolor{white}{\colorbox{clr1hdr}{\parbox{21.0cm}{\centering
  \textbf{Priority Cluster 1 — Sovereign Trusted Infrastructure}
  \quad|\quad Sub-priorities: P1, P2, P3, P8, P9}}}} \\
\hline
\rowcolor{dimHdr}
\textbf{Analytical Dimension} &
\textbf{Phase 1: Foundation \& Compliance (2025--2030)} &
\textbf{Phase 2: Applied \& Industrial R\&D alongside Federation \& Standardisation (2030--2035)} &
\textbf{Phase 3: Autonomous Leadership \& Frontier Innovation (2035--2040)} \\
\hline
\endfirsthead

\hline
\rowcolor{dimHdr}
\textbf{Analytical Dimension} &
\textbf{Phase 1 (2025--2030)} &
\textbf{Phase 2 (2030--2035)} &
\textbf{Phase 3 (2035--2040)} \\
\hline
\endhead

\hline\multicolumn{4}{r}{\footnotesize\textit{Continued\ldots}}\\
\endfoot
\hline
\endlastfoot

% ---- R&I ROW ----
\rowcolor{dimRI}
\cellhead{(1) R\&I Maturity}
\textit{ODP: Applied R\&D}
\textit{[cluster composite; see sub-priority decomposition below]}
&
RISC-V secure element prototypes validated against EUCC evaluation criteria under
EU Chips Act pilot funding; NIST FIPS~203/204/205 algorithms optimised for 8/16-bit
microcontrollers within 32\,KB~RAM constraints, per NIST~IR~8413~performance
benchmarks~\cite{nistir8413}; RF fingerprinting CSI/CFO device authentication
validated in controlled mMTC environments; RISC-V Keystone TEE attestation protocols
tested against heterogeneous edge fleet configurations~\cite{keystone};
soft-state Island Mode credential caching validated in 5G disconnection simulation
environments under SNS~JU governance~\cite{snsju}.
The cluster ODP is Applied R\&D — the lowest common denominator driven by P2, P3, and
P8; P1 and P9 have reached Fragmented Piloting and are not undeployed.
\kpistrip{Transition criterion: at least one functional prototype per sub-priority
validated against the minimum EU compliance baseline (NIS2, CRA, or AI Act as
applicable).}
&
EU-manufactured RISC-V secure elements complete EUCC~`High' evaluation and enter
industrial production; hardware-optimised ML-KEM and ML-DSA validated on EU silicon
for constrained-device targets; RISC-V Keystone TEEs deployed in multi-provider EU
Data Spaces with continuous remote attestation~\cite{keystone}; FHE performance
benchmarks achieve operational viability for defined AI pipeline workloads;
Island Mode swarm coordination protocols validated in multi-operator cross-border
corridor deployments under IPCEI-CIS oversight.
Advanced applied R\&D continues in parallel on zero-energy PhySec, neuromorphic
secure compute, and FHE generalisation — confirming that Phase~2 is not a
purely standardisation phase but a dual-track applied and industrial R\&D stage.
\kpistrip{Transition criterion: production-ready EUCC-evaluated reference
implementation for P1; validated performance benchmarks on EU silicon for P2, P3, P8;
72-hour Island Mode survivability demonstrated operationally for P9.}
&
EU research groups set the international frontier agenda in lightweight PQC,
confidential computing, and physical-layer security without adopting externally
defined research directions. Second-generation PQC hardware co-designed with
EU secure elements; neuromorphic secure edge compute integrated with zero-energy
PhySec; quantum-native TEE architectures for post-CMOS environments;
bio-inspired Island Mode swarm protocols achieving autonomous cyber-physical
recovery at continental scale without any cloud connectivity or central
coordination.
\kpistrip{Transition criterion: the capability is included as a mandatory
baseline requirement in at least one EU-wide standard or procurement framework
with no waiver provisions for new deployments.}
\\
\hline

% ---- POLICY ROW ----
\rowcolor{dimPol}
\cellhead{(2) Policy, Regulatory \& Standards}
\textit{Status: Partially\newline enforced / Standards\newline under development}
&
CRA Articles~10 and~11 full enforcement from October~2027 establishes the mandatory
market condition for hardware security certification investment~\cite{cra}.
NIS2 Article~21(2)(e) cryptographic policy obligations active from October~2024~\cite{nis2}.
NIST FIPS~203/204/205 finalized August~2024 providing algorithm standard
basis~\cite{fips203,fips204,fips205}; ENISA PQC transition guidelines in
publication; EUCC evaluation methodology available but no EU-manufactured secure
element has completed High assurance evaluation; 3GPP~TS~33.501 Release~18
provides 5G security architecture for Island Mode integration~\cite{threeGPP33501}.
DORA Article~9 ICT resilience testing requirements apply to P9~\cite{dora};
AI Act Article~9 risk management obligations apply to TEE-attested autonomous
workloads~\cite{aiact}.
\kpistrip{Critical policy gap: no EU-specific ratified standard for RISC-V
secure element security evaluation existed at the time of writing; this must
be addressed in Phase~1 CEN/CENELEC JTC~13 work programme.}
&
EUCC~`High' certification of the first EU-manufactured RISC-V secure element
issued — the single most critical policy gate condition for this cluster.
CEN/CENELEC hardware security standards for RISC-V secure elements ratified;
ETSI~TS~103~744 hybrid key exchange profiles adopted as mandatory profiles in EU
critical infrastructure procurement~\cite{etsi103744};
EUCS~`High+' certification operational for TEE-backed cloud services~\cite{eucs};
3GPP Release~19/20 specifies PhySec and ISAC security requirements for 6G
non-terrestrial network integration aligned with 3GPP~TR~33.836~\cite{threeGPP33836}.
\kpistrip{Critical policy gate: EUCC `High' certification of the first EU-manufactured
RISC-V secure element is the binding gate condition controlling the Phase~2 SSL
transition for P1, P2, P3, and P8 simultaneously.}
&
EUCC~`High' mandatory baseline for all new EU ICT products in regulated sectors;
EU PQC profiles referenced in ISO/IEC international standards, establishing EU
as a standards-setting jurisdiction rather than an adopter; continuous automated
EUCS/EUCC certification operational at continental scale without manual audit
cycles; AI Act Article~9 conformity assessment for hardware-backed autonomous
security AI fully operational and referenced internationally~\cite{aiact};
DORA Article~11 recovery time objectives achieved autonomously by Island Mode
architecture without programmatic enforcement~\cite{dora}.
\kpistrip{End-state governance criterion: at least one EU standard in this
cluster's capability domain adopted as or formally referenced in an ISO/IEC or
ITU-T international standard.}
\\
\hline

% ---- OPERATIONAL ROW ----
\rowcolor{dimOps}
\cellhead{(3) Operational Deployment \& Adoption}
\textit{Deployment: Research\newline consortia and\newline institutional pilots}
&
Intel TDX and AMD SEV-SNP dominate all production TEE deployments in EU cloud
infrastructure; RISC-V secure elements at prototype stage in 3--5 EU-funded research
consortia; edge nodes universally operate with hard cloud dependencies and no validated
autonomous fallback in production; hybrid PQC pilots operate in cloud research
environments but not in IoT field deployments; Island Mode soft-state protocols
validated in 5G research corridors under SNS~JU programme management~\cite{snsju}.
Deployment actors are limited to large institutional early adopters
(national cybersecurity agencies, telecom operators in EU corridor pilots,
EU-funded research consortia); SME adoption is absent due to tooling immaturity.
\kpistrip{Transition criterion: documented operational deployment across at least
two EU critical infrastructure sectors or Member State jurisdictions with
published findings.}
&
EU RISC-V secure elements deployed in energy and telecommunications critical
infrastructure under IPCEI-CIS programme oversight; hybrid ML-KEM plus ECDH-P384
key encapsulation operational in Industrial IoT corridor deployments across 3+
Member States per ETSI~TS~103~744~\cite{etsi103744};
multi-party confidential computing pipelines active in 3+ EU Data Spaces
with EUCS~`High+' compliance verification; 72-hour Island Mode survivability
validated operationally in 5G cross-border corridor deployments without
WAN connectivity; SME-accessible crypto-agility OTA tooling available via
Digital Europe Programme.
\kpistrip{Transition criterion: deployment across five or more Member States
or three or more EU critical infrastructure sectors, with documented SME
participation in at least one deployment context.}
&
At least 60\% of EU critical infrastructure uses EU-designed EUCC~`High' certified
secure elements as default procurement; all new EU IoT deployments ship with
FIPS~203/204/205-compliant PQC without bespoke integration overhead~\cite{fips203};
FHE operationalised end-to-end for privacy-preserving AI pipelines from IoT to
HPC; bio-inspired Island Mode swarms achieve DORA-mandated recovery objectives
autonomously; Security-by-Design is structurally indistinguishable from product
functionality across the EU ICT market.
\kpistrip{End-state deployment criterion: at least 80\% of new ICT deployments
include this capability as a standard component without dedicated integration effort;
the compliance distinction between ``secure'' and ``standard'' products ceases to
be meaningful.}
\\
\hline

% ---- MARKET/SSL ROW ----
\rowcolor{dimMkt}
\cellhead{(4) Market Uptake \&}\cellhead{Strategic Sovereignty}
\sslbadge{sslredbg}{sslredtx}{SSL: High Foreign Dependency}
&
Non-EU TPMs (Infineon SLB9670, NXP SE050, STMicroelectronics ST33 series) hold
$>$90\% of EU critical infrastructure secure element market under
ISO/IEC~11889~\cite{isoiec11889} and TCG TPM~2.0 specifications~\cite{tgcpm20};
EU Chips Act fabrication pilot lines initiated but first certified secure elements
are 3--5 years from production; no EU-manufactured EUCC~`High' certified component
available in any product category; PQC implementation ecosystem dominated by
US/UK cryptography research outputs; Island Mode tooling absent from EU commercial
product portfolios.
\kpistrip{SSL baseline: High Foreign Dependency across all five sub-priorities.
Market transition criterion: at least one EU-designed solution achieves a verifiable
market entry event (product launch, procurement award, or EUCC/EUCS certification).}
&
\sslbadge{sslamberbg}{sslambtx}{SSL: Emerging EU Ecosystem}\par\smallskip
EU alternatives hold 15--30\% market share in new critical infrastructure secure
element procurement; EU Chips Act fabrication ramping toward 50\% self-sufficiency
target; first commercial EU TEE products (Keystone-based~\cite{keystone})
competing at product level with Intel TDX/AMD SEV-SNP; foreign dependency reduced
from structural to situational in regulated-sector procurement; EU industrial actors
present in sovereign hardware procurement across 5+ Member States.
\kpistrip{Market transition criterion: EU-designed solutions hold at least 30\%
of new critical infrastructure procurement market share, or foreign dependency
formally reduced below the threshold identified at baseline.}
&
\sslbadge{ssltealbg}{ssltealtx}{SSL: Full Sovereign Autonomy}\par\smallskip
EU Chips Act sovereign fabrication meets $\geq$60\% demand for security-critical
silicon across all product categories; EU TEE, secure element, and PQC solutions
commercially exported to partner jurisdictions; zero structural foreign dependency
in new critical infrastructure deployments; EU sets technical terms of competition
in secure silicon globally through EUCC certification requirements and procurement
weight; foreign legacy deployments under active replacement programmes with
binding timelines.
\kpistrip{End-state SSL criterion: EU-designed solutions hold at least 60\%
of new critical infrastructure procurement; foreign legacy under binding
replacement; EU sovereignty is self-sustaining, not programme-managed.}
\\
\hline

% ---- GATE ROW ----
\rowcolor{dimGate}
\textit{\textbf{Critical Enabler}}
\gatelbl{Phase~1 $\rightarrow$ Phase~2}
&
\multicolumn{1}{>{\footnotesize\raggedright\arraybackslash\columncolor{dimGate}}p{5.5cm}|}{%
\textbf{CRA Articles~10 and~11 enforcement (October~2027)~\cite{cra}}
creating mandatory market demand for hardware security certification and
establishing the commercial conditions that justify EU Chips Act fabrication
investment and EUCC~`High' evaluation resource allocation.}
&
\multicolumn{1}{>{\footnotesize\raggedright\arraybackslash\columncolor{dimGate}}p{5.5cm}|}{%
\textbf{EUCC~`High' certification of the first EU-manufactured RISC-V secure
element~\cite{eucc}}
constituting the gate condition that transitions the cluster from research-grade
to market-grade sovereign hardware and unlocks EU Chips Act fabrication scale-up
investment justification.}
&
\multicolumn{1}{>{\footnotesize\raggedright\arraybackslash\columncolor{dimGate}}p{5.5cm}|}{%
\textbf{EU Chips Act sovereign EUV-capable fabrication achieving $\geq$60\%
market self-sufficiency in security-critical silicon}
removing the last structural foreign dependency in the infrastructure trust chain
and enabling Full Sovereign Autonomy as a self-sustaining market condition.}
\\
\hline

% ---- SUB-PRIORITY DECOMPOSITION ----
\rowcolor{dimHdr}
\multicolumn{4}{|>{\footnotesize\itshape\raggedright\arraybackslash}p{21.5cm}|}{%
\textbf{Sub-Priority ODP / SSL Decomposition} —
individual maturity tracking per constituent priority across phases.
Composite cluster ODP and SSL rows above represent the lowest common denominator;
this section preserves intra-cluster maturity heterogeneity.}
\\
\hline

\rowcolor{dimSub}
\textit{P1: Hardware Roots of Trust \& EUCC `High'}
& ODP: Fragmented Piloting \newline SSL: High Foreign Dependency
& ODP: Standardised Integration \newline SSL: Emerging EU Ecosystem
& ODP: Ubiquitous Operation \newline SSL: Full Sovereign Autonomy \\
\hline

\rowcolor{dimSub}
\textit{P2: PQC Agility for Edge/IoT}
& ODP: Applied R\&D \newline SSL: High Foreign Dependency
& ODP: Fragmented Piloting \newline SSL: Emerging EU Ecosystem
& ODP: Ubiquitous Operation \newline SSL: Full Sovereign Autonomy \\
\hline

\rowcolor{dimSub}
\textit{P3: Lightweight Crypto \& Physical-Layer Security}
& ODP: Applied R\&D \newline SSL: High Foreign Dependency
& ODP: Standardised Integration \newline SSL: Emerging EU Ecosystem
& ODP: Ubiquitous Operation \newline SSL: Full Sovereign Autonomy \\
\hline

\rowcolor{dimSub}
\textit{P8: Confidential Computing \& Privacy-Preserving AI}
& ODP: Applied R\&D \newline SSL: High Foreign Dependency
& ODP: Fragmented Piloting \newline SSL: Emerging EU Ecosystem
& ODP: Ubiquitous Operation \newline SSL: Full Sovereign Autonomy \\
\hline

\rowcolor{dimSub}
\textit{P9: Resilient-by-Design \& Autonomous Island Mode}
& ODP: Fragmented Piloting \newline SSL: High Foreign Dependency
& ODP: Standardised Integration \newline SSL: Contained Local Survivability
& ODP: Ubiquitous Operation \newline SSL: Systemic EU Resilience \\
\hline

\end{longtable}

% =========================================================================
% CLUSTER 2 MATRIX
% =========================================================================

\begin{longtable}{%
|>{\footnotesize\raggedright\arraybackslash}p{3.4cm}%
|>{\footnotesize\raggedright\arraybackslash}p{5.5cm}%
|>{\footnotesize\raggedright\arraybackslash}p{5.5cm}%
|>{\footnotesize\raggedright\arraybackslash}p{5.5cm}|}

\caption{Progression Matrix — Priority Cluster~2: Cognitive Security and
  Autonomous Cyber Defence (P6: LLMSecOps, P10: Proactive APT Management,
  AI-as-Target: CL1.2, CL1.3, CL4.7). Assessed across four analytical
  dimensions per the generalised methodology framework
  in this document.}
\label{tab:cluster2_matrix}\\
\hline
\rowcolor{dimHdr}
\multicolumn{4}{|>{\footnotesize\bfseries\centering\arraybackslash}p{21.5cm}|}{%
  \textcolor{white}{\colorbox{clr2hdr}{\parbox{21.0cm}{\centering
  \textbf{Priority Cluster 2 — Cognitive Security and Autonomous Cyber Defence}
  \quad|\quad Sub-priorities: P6, P10, AI-as-Target (CL1.2, CL1.3, CL4.7)}}}} \\
\hline
\rowcolor{dimHdr}
\textbf{Analytical Dimension} &
\textbf{Phase 1: Foundation \& Compliance (2025--2030)} &
\textbf{Phase 2: Applied \& Industrial R\&D alongside Federation \& Standardisation (2030--2035)} &
\textbf{Phase 3: Autonomous Leadership \& Frontier Innovation (2035--2040)} \\
\hline
\endfirsthead

\hline
\rowcolor{dimHdr}
\textbf{Analytical Dimension} &
\textbf{Phase 1 (2025--2030)} &
\textbf{Phase 2 (2030--2035)} &
\textbf{Phase 3 (2035--2040)} \\
\hline
\endhead

\hline\multicolumn{4}{r}{\footnotesize\textit{Continued\ldots}}\\
\endfoot
\hline
\endlastfoot

% ---- R&I ROW ----
\rowcolor{dimRI}
\cellhead{(1) R\&I Maturity}
\textit{ODP: Applied R\&D /\newline Siloed AI Copilots\newline [cluster composite]}
&
NWDAF-integrated ML-based UEBA deployed in 5G standalone pilot environments
per 3GPP~TS~23.288~\cite{threeGPP23288}; EU-sovereign SLMs ($<$10B parameters)
in early training on MITRE ATT\&CK-aligned corpora under GDPR-compliant data
governance~\cite{gdpr}; AI Act Article~9 conformity assessment methodology in
development for autonomous security agents~\cite{aiact}; adversarial robustness
testing frameworks for prompt injection and model evasion under active research;
APT dwell time reduction from 150--200 day baseline targeted via ML-assisted
IoC correlation in pilot enterprise environments.
\kpistrip{Transition criterion: at least one EU-sovereign SLM trained on
MITRE ATT\&CK-aligned EU threat corpora with documented adversarial robustness
validation; at least one NWDAF-integrated UEBA pilot with published dwell-time
reduction evidence.}
&
Edge-native EU-sovereign SLMs certified under AI Act Article~9 conformity
assessment and deployed in critical infrastructure SOC environments;
LLMSecOps agents executing autonomous Zero-Trust reconfigurations per
NIST~SP~800-207~\cite{nistsp800207};
Moving Target Defense (MTD) validated in cross-operator EU corridor simulation
environments per ETSI~GS~ZSM~009 closed-loop automation~\cite{etsizsm009};
RLHF continuous correction pipelines maintaining model accuracy against evolving
APT behavioural profiles; deception technology infrastructure (honeynets,
credential canaries per CL2.6) providing AI training ground truth.
Advanced applied R\&D continues on swarm AI coordination protocols and
self-evolving adversarial robustness frameworks.
\kpistrip{Transition criterion: EU-sovereign SLM deployed operationally in
certified critical infrastructure SOC contexts with documented APT dwell-time
reduction to under 72 hours; adversarial robustness validated under
CEN/CENELEC-standardised framework.}
&
Swarm-based agentic AI coordination protocols active across EU Cloud-Edge
Continuum; self-evolving threat intelligence synthesis via decentralised AI
without central coordination dependency; EU frontier research defines international
agenda in trustworthy autonomous security AI; second-generation adversarial
robustness frameworks address post-quantum AI attack surfaces; validated
APT reconnaissance success rates below 5\% in operational deployment
environments.
\kpistrip{End-state criterion: EU AI security research groups are agenda-setters,
not fast-followers; autonomous APT dwell time below 24 hours as operational norm
across EU critical infrastructure.}
\\
\hline

% ---- POLICY ROW ----
\rowcolor{dimPol}
\cellhead{(2) Policy, Regulatory \& Standards}
\textit{Status: AI Act entering\newline force; conformity\newline assessment capacity\newline building}
&
AI Act Articles~9, 13, and~15 enter into force and apply to autonomous security
containment agents as high-risk AI systems~\cite{aiact};
NIS2 Article~21(2)(b) incident handling procedure obligations create institutional
demand for AI-augmented SOC capabilities~\cite{nis2};
ETSI~GS~ZSM~009 closed-loop automation standards provide the governance framework
for zero-touch security orchestration~\cite{etsizsm009};
MITRE ATT\&CK framework adopted as mandatory behavioural taxonomy for EU threat
intelligence sharing; EU AI Office establishing conformity assessment body capacity
for high-risk AI certification --- this institutional bottleneck represents
the primary Phase~1 policy risk: insufficient conformity assessment capacity
will delay Phase~2 certified operational deployment.
\kpistrip{Critical policy gap: AI Act conformity assessment body capacity for
high-risk autonomous security AI must be established in Phase~1; a 24--36 month
bottleneck in Phase~2 certification is the primary systemic risk for this cluster.}
&
Standardised AI Act Article~9 conformity assessment framework for autonomous
security agents operational, enabling certified deployment without per-case
regulatory approval; ETSI~GS~ZSM~009 formally adopted in EU telco sector
procurement requirements; MITRE ATT\&CK-aligned pan-EU CTI sharing standards
ratified under NIS2 implementing acts~\cite{nis2}; CEN/CENELEC AI security agent
assurance standard published; NIST~SP~800-207 ZTA integration requirements adopted
in EU critical infrastructure security baselines~\cite{nistsp800207}.
\kpistrip{Critical policy gate: standardised AI Act conformity assessment framework
for autonomous security agents is the single gate condition enabling certified
Phase~3 swarm-level operational deployment.}
&
EU AI Act framework for autonomous security AI recognised internationally and
referenced in bilateral digital trade agreements; EU standards for AI security agent
trustworthiness adopted as ISO/IEC international reference;
GDPR Article~22 HITL requirements and AI Act Article~13 explainability
obligations are met by governance-by-default rather than programmatic
enforcement~\cite{gdpr,aiact}; EU as the reference jurisdiction for responsible
autonomous cyber defence AI deployment globally.
\kpistrip{End-state governance criterion: at least one EU AI security standard
adopted as or referenced in an ISO/IEC or ITU-T international standard; EU
regulatory framework cited in at least two bilateral digital trade agreements.}
\\
\hline

% ---- OPERATIONAL ROW ----
\rowcolor{dimOps}
\cellhead{(3) Operational Deployment \& Adoption}
\textit{Deployment: US-domiciled\newline SIEM dominant;\newline AI copilots in\newline pilot SOCs}
&
LLMs operating as analyst copilots in EU enterprise SOC environments with human
analysts retaining all containment authority; NWDAF telemetry routed to centralised
AI analytics in 5G standalone deployments of 3+ EU operators; US-domiciled SIEM
platforms (Splunk/Cisco, IBM QRadar) augmented with ML correlation layers;
adversarial robustness validation confined to research environments with no
operational certification framework in force; no EU-sovereign foundation model for
security operations in production deployment; APT containment remains manual
with AI-assisted alert triage only.
\kpistrip{Transition criterion: documented AI-augmented SOC deployment in at least
two EU critical infrastructure sectors with published dwell-time reduction evidence;
at least one EU-sovereign SLM pilot in an operational security context.}
&
EU-sovereign SLMs deployed in critical infrastructure SOC environments across
5+ Member States; autonomous Zero-Trust reconfigurations operational in certified
EU 5G core deployments; LLMSecOps agents ingesting CTI and tracking lateral
movement in real time against MITRE ATT\&CK behavioural frameworks; APT dwell
time reduced to under 72 hours in EU critical infrastructure under autonomous
agent operation; purple-team automation (CL1.4) conducting continuous adversarial
validation against live staging environments; deception technology (CL2.6)
generating high-fidelity APT ground truth for continuous agent retraining.
\kpistrip{Transition criterion: EU-sovereign SLM deployed operationally in 5+
Member States; autonomous containment certified under AI Act Article~9 for at
least one critical infrastructure sector.}
&
Swarm AI defends EU Cloud-Edge Continuum autonomously with APT reconnaissance
success rates below 5\% in operational environments; Moving Target Defense
reshuffles network topology, IP addressing, and software diversity at continental
scale per ETSI~GS~ZSM~009 zero-touch automation~\cite{etsizsm009};
deception infrastructure generates real-time adversarial ground truth continent-wide;
autonomous agents meet AI Act Article~13 explainability requirements without
human intervention for standard containment decisions; zero-touch SOC operation
is the operational baseline for all EU critical infrastructure.
\kpistrip{End-state deployment criterion: zero-touch security orchestration is the
operational default; manual analyst intervention reserved for novel threat
escalation only, not routine containment.}
\\
\hline

% ---- MARKET/SSL ROW ----
\rowcolor{dimMkt}
\cellhead{(4) Market Uptake \&}\cellhead{Strategic Sovereignty}
\sslbadge{sslredbg}{sslredtx}{SSL: High Reliance on US Foundation Models}
&
US-domiciled SIEM platforms (Splunk/Cisco, IBM QRadar) dominate EU enterprise
SOC deployments; EU organisations structurally dependent on non-EU threat
intelligence feeds (Mandiant/Google, CrowdStrike, Microsoft Sentinel) for APT
indicator attribution; no EU-sovereign LLM for security operations at
production scale; EuroHPC AI Factory compute capacity insufficient for training
frontier security foundation models within Phase~1 timeframe, constraining EU
models to SLM/LCM scale ($<$10B parameters).
\kpistrip{SSL baseline: High Reliance on US Foundation Models / High Dependency
on Global Threat Feeds. Market transition criterion: at least one EU-sovereign
SLM achieves a production deployment event in an EU critical infrastructure SOC.}
&
\sslbadge{sslamberbg}{sslambtx}{SSL: Sovereign EU Edge LLMs / SLMs Emerging}\par\smallskip
EU-sovereign SLMs deployed in $\geq$30\% of critical infrastructure SOC
environments; US-domiciled SIEM dependency reduced to legacy deployments in
regulated sectors under active replacement; EU CTI sharing infrastructure
(CONCORDIA successor~\cite{concordia}) operational as sovereign alternative
to US-domiciled threat intelligence feeds; EuroHPC AI Factory compute
sufficient for training EU security SLMs at operational scale; first EU AI
security products exported to partner jurisdictions.
\kpistrip{Market transition criterion: EU-sovereign SLMs hold at least 30\%
of EU critical infrastructure SOC deployments; measurable reduction in
US-domiciled SIEM dependency for new procurement in regulated sectors.}
&
\sslbadge{ssltealbg}{ssltealtx}{SSL: Global Leadership in Trustworthy AI Security}\par\smallskip
EU-sovereign SLMs cover $\geq$50\% of EU critical infrastructure SOC
deployments; US-domiciled SIEM dependency eliminated for new deployments in
regulated sectors; EU AI security technology exported to and adopted by partner
jurisdictions as governance and technology reference; EU as global standard-setter
for trustworthy autonomous cyber defence; continental APT dwell time from
150--200 day baseline reduced to under 24 hours as operational norm.
\kpistrip{End-state SSL criterion: EU AI security models adopted as reference
implementations by partner jurisdictions; EU is a net exporter of AI security
governance standards, not an importer.}
\\
\hline

% ---- GATE ROW ----
\rowcolor{dimGate}
\textit{\textbf{Critical Enabler}}
\gatelbl{Phase transition}
&
\multicolumn{1}{>{\footnotesize\raggedright\arraybackslash\columncolor{dimGate}}p{5.5cm}|}{%
\textbf{AI Act Articles~9 and~15 enforcement (2025--2026)~\cite{aiact}}
establishing mandatory risk management and adversarial robustness requirements
for autonomous security AI, creating the regulatory framework for certified
Phase~2 deployment and incentivising EU-sovereign model development as the
compliance-safe alternative to US-domiciled foundation models.}
&
\multicolumn{1}{>{\footnotesize\raggedright\arraybackslash\columncolor{dimGate}}p{5.5cm}|}{%
\textbf{Standardised AI Act conformity assessment framework for autonomous
security agents}
enabling certified operational deployment at scale without per-case regulatory
approval, removing the certification bottleneck that otherwise delays Phase~3
swarm-level operational autonomy across all three clusters.}
&
\multicolumn{1}{>{\footnotesize\raggedright\arraybackslash\columncolor{dimGate}}p{5.5cm}|}{%
\textbf{EU-certified self-evolving AI swarm protocols achieving full operational
deployment under ETSI~GS~ZSM~009 zero-touch automation governance~\cite{etsizsm009}}
with pan-EU CTI-sharing infrastructure providing sovereign adversarial intelligence
at continental scale.}
\\
\hline

% ---- SUB-PRIORITY DECOMPOSITION ----
\rowcolor{dimHdr}
\multicolumn{4}{|>{\footnotesize\itshape\raggedright\arraybackslash}p{21.5cm}|}{%
\textbf{Sub-Priority ODP / SSL Decomposition} —
individual maturity tracking per constituent priority across phases.}
\\
\hline

\rowcolor{dimSub}
\textit{P6: Cognitive SecOps (LLMSecOps) \& Automated Threat Intelligence}
& ODP: Siloed AI Copilots \newline SSL: High Reliance on US Foundation Models
& ODP: Federated Edge Agents \newline SSL: Sovereign EU Edge LLMs Emerging
& ODP: Ubiquitous AI Immune System \newline SSL: Global Leadership in Trustworthy AI \\
\hline

\rowcolor{dimSub}
\textit{P10: Proactive APT Management}
& ODP: Multi-Layered Alert Correlation \& Hunting \newline SSL: High Dependency on Global Threat Feeds
& ODP: Automated Triage \& LLMSecOps Containment \newline SSL: Sovereign Pan-EU CTI \& Behavioural Analytics
& ODP: Ubiquitous Moving Target Defense (MTD) \newline SSL: Full Autonomous Counter-Adversarial Resilience \\
\hline

\rowcolor{dimSub}
\textit{AI-as-Target (CL1.2, CL1.3, CL4.7): Model Integrity, Runtime SBOMs, AI Trustworthiness}
& ODP: Applied R\&D \newline SSL: High Foreign Dependency
& ODP: Standardised Integration \newline SSL: Emerging EU Ecosystem
& ODP: Ubiquitous Operation \newline SSL: Full Sovereign Autonomy \\
\hline

\end{longtable}

% =========================================================================
% CLUSTER 3 MATRIX
% =========================================================================

\begin{longtable}{%
|>{\footnotesize\raggedright\arraybackslash}p{3.4cm}%
|>{\footnotesize\raggedright\arraybackslash}p{5.5cm}%
|>{\footnotesize\raggedright\arraybackslash}p{5.5cm}%
|>{\footnotesize\raggedright\arraybackslash}p{5.5cm}|}

\caption{Progression Matrix — Priority Cluster~3: Federated Operational
  Sovereignty (P4: Security Interoperability via Simpl, P5: Supply Chain
  Verification, P7: Pan-European NaaS Federation). Assessed across four
  analytical dimensions per the generalised methodology framework
  in this document.}
\label{tab:cluster3_matrix}\\
\hline
\rowcolor{dimHdr}
\multicolumn{4}{|>{\footnotesize\bfseries\centering\arraybackslash}p{21.5cm}|}{%
  \textcolor{white}{\colorbox{clr3hdr}{\parbox{21.0cm}{\centering
  \textbf{Priority Cluster 3 — Federated Operational Sovereignty}
  \quad|\quad Sub-priorities: P4, P5, P7}}}} \\
\hline
\rowcolor{dimHdr}
\textbf{Analytical Dimension} &
\textbf{Phase 1: Foundation \& Compliance (2025--2030)} &
\textbf{Phase 2: Applied \& Industrial R\&D alongside Federation \& Standardisation (2030--2035)} &
\textbf{Phase 3: Autonomous Leadership \& Frontier Innovation (2035--2040)} \\
\hline
\endfirsthead

\hline
\rowcolor{dimHdr}
\textbf{Analytical Dimension} &
\textbf{Phase 1 (2025--2030)} &
\textbf{Phase 2 (2030--2035)} &
\textbf{Phase 3 (2035--2040)} \\
\hline
\endhead

\hline\multicolumn{4}{r}{\footnotesize\textit{Continued\ldots}}\\
\endfoot
\hline
\endlastfoot

% ---- R&I ROW ----
\rowcolor{dimRI}
\cellhead{(1) R\&I Maturity}
\textit{ODP: Fragmented\newline Piloting [P4, P7];\newline Supply Chain\newline Disclosure [P5]}
&
Initial Simpl security MIM pilots in 3 EU Data Spaces (Manufacturing, Health, Energy);
CRA-baseline SBOM tooling deployed by early adopters in Open RAN supply chains;
CAMARA API harmonisation tested in 2--3 bilateral 5G corridor agreements;
automated compliance negotiator prototypes validated in controlled IPCEI-CIS
pilot environments; semantic security ontology (STIX~2.1/TAXII) integration
for federated SOC interoperability under active R\&I; ZKP computational overhead
benchmarking for supply chain attestation use cases confirming batch verification
viability.
\kpistrip{Transition criterion: at least one functional Simpl MIM pilot
demonstrated in a named Data Space with automated cross-provider security
policy enforcement verified by an independent audit.}
&
Automated Simpl MIM policy translation operational across $\geq$50\% of active
EU Data Spaces; Runtime SBOMs with cryptographic proof of execution verifying
containerised components dynamically across multi-provider continuum nodes;
ZKP benchmarks achieving batch verification viability for supply chain attestation;
CAMARA-based cross-border edge compute roaming active across 5+ EU operator
pairs; Data Visa mechanisms enabling AI model migration without manual legal
review; inter-operator east-west interface standards ratified and operational
in IPCEI-CIS corridor deployments.
Advanced applied R\&D: next-generation semantic interoperability for
post-quantum supply chains; federated digital twin infrastructure (CL3.5)
for pre-production security validation.
\kpistrip{Transition criterion: Runtime SBOM with cryptographic execution proof
operational in at least one Open RAN deployment; CAMARA cross-border roaming
active across 5+ operator pairs with documented SLA compliance.}
&
Provider-agnostic EU trust fabric operational continent-wide with real-time
automated security orchestration; ZKP-verified supply chain attestation
eliminates unverified components from EU critical infrastructure at continental
scale; unified EU NaaS platform legally operable as one data jurisdiction;
EU supply chain governance frameworks adopted as international reference models;
federated cyber range infrastructure (CL3.5) providing pre-production validation
for all cross-border security configurations before continental deployment.
\kpistrip{End-state criterion: EU supply chain governance and NaaS federation
frameworks referenced in at least two international standards or bilateral
digital trade agreements.}
\\
\hline

% ---- POLICY ROW ----
\rowcolor{dimPol}
\cellhead{(2) Policy, Regulatory \& Standards}
\textit{Status: CRA/NIS2/DORA\newline enforcement creating\newline supply chain and\newline interop mandates}
&
CRA Annex~I Part~II SBOM disclosure obligations create the mandatory supply chain
transparency baseline~\cite{cra}; NIS2 Article~21(2)(h) supply chain security
measures mandatory for essential entities~\cite{nis2};
DORA Article~28 ICT third-party risk management requirements active for
financial sector~\cite{dora}; Data Act interoperability obligations entering
into force; Simpl-Open official launch with initial security MIM definitions
establishing the technical baseline for CEN/CENELEC standardisation;
GSMA CAMARA API standards reaching v1.x publication milestone.
\kpistrip{Critical policy gap: Data Act implementing acts specifying interoperability
technical requirements for security interfaces must be finalised within Phase~1
to provide standards basis; delayed publication creates a compliance gap for
cross-border Data Space deployments.}
&
Security MIMs adopted as CEN/CENELEC standards with mandatory reference in EU
Cloud Rulebook — the single binding gate condition for non-voluntary hyperscaler
compliance; DORA Article~28 continuous third-party monitoring requirements creating
procurement pull for runtime SBOM verification tooling; Data Act full enforcement
driving interoperability adoption beyond early adopters; EU NaaS federation legal
framework (Data Visa mechanisms) in active legislative development;
EUCS automated certification operational for Simpl-federated cloud
services~\cite{eucs}.
\kpistrip{Critical policy gate: security MIMs adopted as CEN/CENELEC standards
with mandatory Cloud Rulebook reference — this is the single gate condition
controlling the P4 and P7 SSL transition from Emerging EU Ecosystem to Full
Sovereign Autonomy.}
&
``Schengen for Cloud'' legal framework enacted, resolving US CLOUD Act
extra-territorial jurisdiction conflict and enabling EU NaaS to operate as a
single data jurisdiction; EU supply chain governance model referenced
internationally; EUCS automated certification adopted in bilateral trade
agreements as mutual recognition framework; EU Cloud Rulebook MIM requirements
referenced in ISO/IEC service interoperability standards; ZKP supply chain
verification standard adopted as international reference.
\kpistrip{End-state governance criterion: Schengen for Cloud enacted; at least
two bilateral mutual recognition agreements referencing EU supply chain
certification; EU interoperability standards referenced in ISO/IEC.}
\\
\hline

% ---- OPERATIONAL ROW ----
\rowcolor{dimOps}
\cellhead{(3) Operational Deployment \& Adoption}
\textit{Deployment: Hyperscaler\newline API dominance;\newline bilateral pilots\newline only}
&
Static CRA-baseline SBOM declarations submitted as point-in-time documents without
runtime verification; CAMARA APIs tested in 2--3 bilateral 5G corridor agreements
between individual EU operators; proprietary hyperscaler APIs (AWS, Azure, GCP)
dominate all cross-provider integration; Simpl security MIM pilots deployed in
proof-of-concept contexts in Manufacturing and Health Data Spaces;
supply chain visibility confined to tier-1 vendor declarations without
cryptographic verification; no operational pan-EU NaaS platform.
Deployment actors: EU-funded research consortia and telecom operators in
designated corridor pilots; SME participation absent due to compliance cost.
\kpistrip{Transition criterion: operational Simpl MIM deployment in at least
two EU Data Spaces verified by independent audit; at least three EU telecom
operators in bilateral CAMARA-based cross-border roaming arrangement.}
&
Automated Simpl MIM policy translation operational across $\geq$50\% of active EU
Data Spaces with cryptographically verifiable enforcement backed by Cluster~1
TEE attestation; Runtime SBOMs with execution proofs deployed in Open RAN supply
chains of 3+ EU operators; CAMARA-standardised cross-border edge compute roaming
active across 5+ EU operator pairs; supply chain quarantine (eBPF-isolated slices)
operational for unverified components; SME-accessible SBOM and MIM compliance
tooling available via Digital Europe Programme support schemes.
\kpistrip{Transition criterion: deployment across 5+ Member States or 3+ critical
infrastructure sectors; documented SME participation; supply chain quarantine
operational in at least one production environment.}
&
All EU Common Data Spaces operating on Simpl-based security interoperability
with automated MIM compliance verification; ZKP-verified supply chain attestation
integrated into all EU critical infrastructure procurement workflows without human
sign-off; unified NaaS platform enabling EU operators to compete at hyperscaler
scale within the EU regulated market; supply chain quarantine automated with
zero-human intervention; federated cyber range (CL3.5) providing pre-production
validation for cross-border configurations before continental deployment.
\kpistrip{End-state deployment criterion: security interoperability is the
unconditional default; no EU critical infrastructure procurement proceeds without
Simpl MIM compliance verification; supply chain clarity is structurally
guaranteed, not audited.}
\\
\hline

% ---- MARKET/SSL ROW ----
\rowcolor{dimMkt}
\cellhead{(4) Market Uptake \&}\cellhead{Strategic Sovereignty}
\sslbadge{sslamberbg}{sslambtx}{SSL: Emerging EU Ecosystem [P4] / Hyperscaler Dominance [P7]}
&
AWS, Azure, and GCP collectively control approximately 65\% of EU cloud market
with proprietary API lock-in; no operational pan-EU NaaS platform; SBOM ecosystem
fragmented between US-tooling vendors (Anchore, Snyk) and immature EU alternatives;
SME participation constrained by compliance cost for CRA-mandated supply chain
verification tooling~\cite{cra}; US CLOUD Act extra-territorial jurisdiction
remains a structural sovereignty risk for all EU data processed on
non-EU-domiciled infrastructure.
\kpistrip{SSL baseline: Emerging EU Ecosystem (P4); Hyperscaler Platform Dominance
(P7); High reliance on fragmented non-EU supply chain auditing (P5). Market
transition criterion: at least one EU MIM-compliant cloud service achieves a
verifiable procurement award in a regulated-sector Data Space.}
&
\sslbadge{ssltransbg}{ssltrantx}{SSL: Growing Pan-EU Sovereignty}\par\smallskip
EU-based cloud providers holding measurably growing share in regulated-sector
procurement via Simpl MIM compliance; hyperscaler API lock-in reduced to legacy
deployments as MIM-compliant alternatives gain procurement weight in 5+ Member
States; EU NaaS consortium beginning to compete with US hyperscaler scale in
defined vertical sectors; supply chain dependency on non-EU tier-1 vendors
reduced in Open RAN deployments; ZKP-verified supply chain attestation enabling
EU actors to demonstrate compliance superiority over foreign self-certification
declarations.
\kpistrip{Market transition criterion: EU MIM-compliant cloud services hold
at least 30\% of new regulated-sector Data Space procurement; measurable reduction
in hyperscaler API lock-in for new critical infrastructure deployments.}
&
\sslbadge{ssltealbg}{ssltealtx}{SSL: Full Federated Sovereign Autonomy}\par\smallskip
EU cloud and edge providers achieve hyperscaler-comparable scale within EU
regulated market through NaaS federation; US CLOUD Act jurisdiction exposure
eliminated for EU data processed on federated sovereign infrastructure;
supply chain transparency sufficient to exclude non-compliant components without
bilateral negotiation; EU interoperability and supply chain standards exported
to and adopted by partner jurisdictions; structural foreign dependency in cloud,
edge, and supply chain domains reduced to managed residual legacy under active
replacement programmes.
\kpistrip{End-state SSL criterion: EU NaaS achieves hyperscaler-comparable scale
in EU regulated market; Schengen for Cloud eliminates CLOUD Act exposure for
new deployments; EU is net exporter of supply chain governance standards.}
\\
\hline

% ---- GATE ROW ----
\rowcolor{dimGate}
\textit{\textbf{Critical Enabler}}
\gatelbl{Phase transition}
&
\multicolumn{1}{>{\footnotesize\raggedright\arraybackslash\columncolor{dimGate}}p{5.5cm}|}{%
\textbf{Simpl-Open official launch with security MIM definitions published as
open standards~\cite{cloud_edge_roadmap}}
establishing the technical baseline for cross-provider security policy
automation and providing the specification reference that makes
CEN/CENELEC standardisation in Phase~2 possible.}
&
\multicolumn{1}{>{\footnotesize\raggedright\arraybackslash\columncolor{dimGate}}p{5.5cm}|}{%
\textbf{Security MIMs adopted as CEN/CENELEC standards with mandatory reference
in EU Cloud Rulebook}
creating the regulatory basis for non-voluntary hyperscaler compliance and
transforming voluntary interoperability into a structural market requirement
enforceable against all cloud service providers operating in the EU market.}
&
\multicolumn{1}{>{\footnotesize\raggedright\arraybackslash\columncolor{dimGate}}p{5.5cm}|}{%
\textbf{``Schengen for Cloud'' legal framework enacted}
resolving US CLOUD Act extra-territorial jurisdiction conflict and enabling
the EU NaaS platform to operate legally as a single data jurisdiction,
achieving the commercial scale conditions for Full Federated Sovereign Autonomy.}
\\
\hline

% ---- SUB-PRIORITY DECOMPOSITION ----
\rowcolor{dimHdr}
\multicolumn{4}{|>{\footnotesize\itshape\raggedright\arraybackslash}p{21.5cm}|}{%
\textbf{Sub-Priority ODP / SSL Decomposition} —
individual maturity tracking per constituent priority across phases.}
\\
\hline

\rowcolor{dimSub}
\textit{P4: Security Interoperability via Simpl Middleware}
& ODP: Fragmented Piloting \newline SSL: Emerging EU Ecosystem
& ODP: Standardised Integration \newline SSL: Full Sovereign Autonomy
& ODP: Ubiquitous Operation \newline SSL: Full Sovereign Autonomy \\
\hline

\rowcolor{dimSub}
\textit{P5: Supply Chain Verification \& Secured Infrastructure}
& ODP: Full Supply Chain Disclosure \newline SSL: High Reliance on Fragmented Non-EU Auditing
& ODP: Dynamic Runtime Verification \newline SSL: Emerging EU Supply Chain Tracking Frameworks
& ODP: ZKP-Verified Attestation \& Full Restriction \newline SSL: Full Sovereign Supply Chain Control \\
\hline

\rowcolor{dimSub}
\textit{P7: Pan-European NaaS Federation}
& ODP: Bilateral Interconnection \newline SSL: Hyperscaler Platform Dominance
& ODP: Federated Edge Roaming \newline SSL: Pan-EU Consortium Competes
& ODP: True Digital Single Market \newline SSL: European Scale Parity Achieved \\
\hline

\end{longtable}

\end{landscape}

\end{document}
